\section{匿名化为什么不是隐私}
\label{sec:17-Constrained-Guarantees/01-Differential-Privacy/01}

一家医院要把住院记录交给外部研究团队。法务的要求写得很清楚：姓名、身份证号、住址、电话全部删掉。工程师照做了，剩下的字段是出生日期、五位邮编、性别、入院日期和诊断编码——看上去没有一列能指向具体的人。

几个月后，一位研究者拿这份数据做了一件没人预料的事。他从公开渠道买到了同一个州的选民登记表，那份表里有姓名、地址、出生日期、邮编和性别。两张表按``出生日期 + 邮编 + 性别''这三列做连接，相当一部分记录只匹配上唯一一个人。诊断编码于是有了名字。

这件事的教训常被总结成``脱敏做得不够彻底''，但这个总结是错的，而且错得很有代表性。删除标识列的思路默认了一件事：隐私风险附着在某些特定字段上，把那些字段拿掉，风险就随之消失。真实情况是，风险附着在\textbf{字段的组合}上。出生日期单独看有几万种取值，邮编单独看有几万种，性别只有两种；三者组合起来的取值空间远大于一个邮编区里的居民数，于是绝大多数人在这个组合上是唯一的。没有任何一列是标识符，它们合起来是。

\subsection{删列这条路走不通，原因不在工程}

发现问题之后，自然的补救是继续删：把出生日期粗化成出生年份，把五位邮编截成三位。这确实压低了唯一性。但这条路有一个结构性的问题——它没有终点。

粗化到年份和三位邮编之后，仍然可能有别的组合能重识别：入院日期加诊断编码，或者多次住院构成的时间序列。每堵住一处，攻击者换一处。而防守方在设计粗化方案时，必须预先猜到攻击者手上有哪些外部数据，这个猜测无法被验证——今天不存在的公开数据集，明天可能存在。安全性建立在``攻击者没有某份数据''这个假设上，而这个假设会随时间单调地变坏。

另一头同样麻烦。粗化到只剩省份和年龄段的数据集确实很难重识别，也几乎无法支撑任何研究。数据里剩下的信息量和它能被重识别的可能性来自同一个源头：如果记录之间还有可分辨的差异，这些差异就还能拿来匹配；如果差异被抹平到无法匹配，研究者要找的关联也一并被抹平了。这不是实现得不够好，是这条路本身在两端各有一个不能同时满足的要求。

\subsection{换一个问法}

死路通常意味着问题问错了。``输出里还残留多少关于个人的信息''这个问题之所以没法回答，是因为``关于个人的信息''没有边界——任何一个统计量都或多或少携带每个人的痕迹，无法把它切成``属于张三的部分''和``不属于张三的部分''。

换成这样问：假设张三的那条记录根本不在数据集里，发布出去的东西会有什么不同？如果``在''与``不在''产生的输出几乎无法区分，那么任何看到输出的人，无论他手上还有多少外部数据，都无法从输出中判断张三是否参与过——他对张三的了解不会因为这次发布而增加。

这个问法有三个好处。它把保护对象从``某个字段''换成了``某个人的参与与否''，后者是可以精确定义的；它把``几乎无法区分''这件事交给概率来度量，于是可以量化；它不需要假设攻击者知道什么，因为``两种情况下的输出分布几乎相同''这个陈述与攻击者的背景知识无关。

要让这个想法成立，发布过程必须是随机的。确定性的机制在 \(D\) 和 \(D'\) 上给出两个确定的数，只要这两个数不同，区分就是完美的。所以差分隐私从一开始就假定输出是随机变量，讨论的是它的分布。

\begin{definition}[相邻数据集与 \(\varepsilon\)-差分隐私]
称两个数据集 \(D,D'\) 相邻（记 \(D\sim D'\)），若它们至多在一条记录上不同。随机机制 \(\mathcal M\) 满足 \(\varepsilon\)-差分隐私，若对\textbf{任意}一对相邻数据集 \(D\sim D'\) 与\textbf{任意}可测输出集合 \(S\)，
\[
P\big(\mathcal M(D)\in S\big)\;\le\;e^{\varepsilon}\,P\big(\mathcal M(D')\in S\big).
\]
\end{definition}

把 \(D\) 与 \(D'\) 互换角色，同一条不等式给出反向的界，于是两个概率之比被夹在 \(e^{-\varepsilon}\) 与 \(e^{\varepsilon}\) 之间。\(\varepsilon\) 越小，两个分布越难分辨；\(\varepsilon=0\) 意味着输出与数据完全无关，也就完全没用。

\subsection{强度全部来自那两个``任意''}

这个定义的力量不在指数函数上，而在两个全称量词上，值得逐个拆开看。

``对任意相邻数据集''管的是最坏情况。它没有说``对大多数数据集''，也没有说``对典型的记录''。哪怕数据集里有一条极端离群的记录——某个罕见病例、某笔异常大额交易——保证对它同样成立。一旦把这个量词弱化成``对随机抽取的相邻数据对以高概率成立''，最需要保护的那些人恰好落在例外里。

``对任意输出集合''管的是攻击者的自由度。攻击者不必事先声明他要看什么，他可以在看到输出之后，构造任何一个他想要的事件 \(S\)，不等式仍然成立。如果只要求对某些特定的 \(S\)（比如``输出大于某个阈值''）成立，攻击者只需换一个更精巧的 \(S\) 就能突破。

这正是\hyperref[sec:01-Mathematical-Language/01-Logic/03]{``量词与否定''}反复强调的那件事：一句数学陈述的实际内容由量词的顺序和作用范围决定，而不是由它读起来的语气决定。把这个定义否定一次会看得更清楚——违反 \(\varepsilon\)-差分隐私意味着\textbf{存在}一对相邻数据集和\textbf{存在}一个输出集合，使不等式失败。一个反例就足以推翻整个保证，这是全称命题的标准形状，也是\hyperref[sec:01-Mathematical-Language/03-Proof-Techniques/03]{``存在性与唯一性''}里那套对偶关系的直接应用。

还有一点更容易被误用：这是\textbf{关于机制的性质，不是关于某次输出的性质}。拿到一个具体的发布结果，问``这个结果隐私吗''是没有意义的，就像拿到一个具体的置信区间问``这次真值在里面吗''一样没有意义。\hyperref[sec:11-Mathematical-Statistics/04-Interval-Estimation/01]{``置信度属于程序，而非这次参数''}讲的是同一种判断：\(95\%\) 描述的是造区间那套程序在重复使用下的表现，\(\varepsilon\) 描述的是产生输出那套程序在所有可能输入上的表现。两者都是对程序的断言，一旦落到单次实例上就失去了内容。这个区分不是学究气——它决定了一件事能不能事后补救。程序的性质必须在发布之前论证清楚，发布之后再去检查输出``看起来安不安全''，检查的是一个没有定义的量。

\begin{center}
\begin{tikzpicture}[x=1.35cm,y=1cm,>=stealth,font=\small]
  % ---- 上：两个输出分布 ----
  \draw[->,black!60] (-0.7,0) -- (6.95,0) node[right,font=\small] {输出};
  \draw[thick,blue!65] plot[domain=-0.6:6.6,samples=100]
    (\x,{1.5*exp(-(\x-2.6)*(\x-2.6)/1.6)});
  \draw[thick,orange!85!black] plot[domain=-0.6:6.6,samples=100]
    (\x,{1.5*exp(-(\x-3.4)*(\x-3.4)/1.6)});
  \node[blue!65,font=\footnotesize] at (1.15,1.55) {\(\mathcal M(D)\)};
  \node[orange!85!black,font=\footnotesize] at (4.9,1.55) {\(\mathcal M(D')\)};
  \node[align=left,anchor=west,font=\footnotesize,black!70] at (-0.6,2.25)
    {\(D\) 与 \(D'\) 只差一条记录};

  % ---- 下：对数密度比 ----
  \begin{scope}[shift={(0,-3.6)}]
    \fill[blue!8] (-0.6,-0.62) rectangle (6.6,0.62);
    \fill[orange!22] (-0.6,-0.88) rectangle (0,0.88);
    \fill[orange!22] (6,-0.88) rectangle (6.6,0.88);
    \draw[black!40,dashed] (-0.6,0.62) -- (6.6,0.62);
    \draw[black!40,dashed] (-0.6,-0.62) -- (6.6,-0.62);
    \draw[->,black!60] (-0.6,0) -- (6.95,0) node[right,font=\small] {输出};
    \draw[very thick,black!80] (-0.6,0.75) -- (6.6,-0.75);
    \node[left,font=\footnotesize] at (-0.66,0.62) {\(+\varepsilon\)};
    \node[left,font=\footnotesize] at (-0.66,-0.62) {\(-\varepsilon\)};
    \node[align=left,anchor=west,font=\footnotesize,black!70] at (-0.6,1.55)
      {对数密度比 \(\log\dfrac{p_D(t)}{p_{D'}(t)}\)};
    \node[align=center,font=\footnotesize,orange!75!black] at (6.3,1.5)
      {\(\delta\) 兜住\\的例外};
  \end{scope}
\end{tikzpicture}
\end{center}

\emph{要记住的是下半张图：差分隐私管的不是``输出里有没有姓名''，而是两条输出分布的比值在整条实轴上被夹在一条固定的带子里；带子的宽度就是 \(\varepsilon\)，跑出带子的那一小块要由 \(\delta\) 单独付账。}

\subsection{\texorpdfstring{松一档：\((\varepsilon,\delta)\) 版本}{松一档：(eps, delta) 版本}}

图里那条对数比值直线在两端跑出了带子，这不是画图时的疏忽。很多有用的机制——包括下一节要讲的 Gaussian 机制——在尾部就是这样：绝大多数输出上比值受控，极少数输出上完全失控。严格的 \(\varepsilon\)-差分隐私容不下这种情况，于是有了松一档的版本。

\begin{definition}[\((\varepsilon,\delta)\)-差分隐私]
机制 \(\mathcal M\) 满足 \((\varepsilon,\delta)\)-差分隐私，若对任意 \(D\sim D'\) 与任意 \(S\)，
\[
P\big(\mathcal M(D)\in S\big)\;\le\;e^{\varepsilon}\,P\big(\mathcal M(D')\in S\big)+\delta.
\]
\end{definition}

多出来的 \(\delta\) 是一个加性余项，它的含义是``允许有概率不超过 \(\delta\) 的一小块区域，在那里保证彻底不成立''。注意``彻底''这个词：在例外区域里，定义不再提供任何限制，输出可以把整条记录原样泄露出去。

这决定了 \(\delta\) 该怎么取。考虑一个荒唐但完全合法的机制：以概率 \(\delta\) 直接公开随机挑出的若干条完整记录，否则输出一个与数据无关的常数。它满足 \((0,\delta)\)-差分隐私，因为例外区域的概率恰好是 \(\delta\)。如果有人把 \(\delta\) 设成 \(10^{-3}\)，而数据集有 \(n=100\) 人，那么这个机制每千次调用就会完整泄露一次——参数表上看却是漂亮的 \(\varepsilon=0\)。

\begin{remark}
实践中要求 \(\delta\ll 1/n\)，常见取法是 \(\delta=n^{-2}\) 或 \(\delta=n^{-1.5}\)。理由就是上面那个构造：只要 \(\delta\) 达到 \(1/n\) 的量级，``随机泄露若干条完整记录''这类机制就能挤进定义，而定义本身没有办法把它们排除出去。这条经验规则不是定理，它是针对一族已知反例设的门槛；换一族反例，门槛还得重新算。
\end{remark}

\subsection{条件对账：这个定义保护什么，不保护什么}

定义写清楚之后，要逐条核对它承诺了什么。

\textbf{相邻关系的定义决定了保护单位。}上面用的是``至多一条记录不同''，这是替换型相邻。另一种常见写法是``一个数据集比另一个多一条记录''，即增删型相邻。两者给出的 \(\varepsilon\) 相差一个大约两倍的因子，在比较不同工作的参数时必须先对齐这一条。更要紧的是：如果一个人在数据集里有多条记录——同一个患者多次住院——那么``差一条记录''保护的是单次住院，不是这个人。保护到人的层面需要把相邻关系定义成``差一个人的全部记录''，敏感度随之上升，同样的 \(\varepsilon\) 要加更多噪声。这不是理论上的细枝末节，它是实际部署中最容易搞错的一处。

\textbf{\(\varepsilon\) 的语义是后验比值的上界。}设攻击者对``张三是否在数据集中''有一个先验，看到输出 \(t\) 之后按 Bayes 规则更新（见\hyperref[sec:05-Probability/02-Conditional-Probability/02]{``全概率公式与 Bayes 更新''}）。由定义，似然之比被夹在 \(e^{\pm\varepsilon}\) 之间，于是后验几率与先验几率之比也被夹在同一区间里。\(\varepsilon=0.1\) 意味着几率至多变化约 \(10\%\)，\(\varepsilon=5\) 意味着几率可以变化一百多倍——后者在论文里仍被称作``满足差分隐私''，但它给出的实际保证已经很弱。报告 \(\varepsilon\) 而不解释这个换算，读者容易把它当成一个越小越好的调参数字，而看不出它在哪个量级上开始失去意义。

\textbf{它不保护群体层面的结论。}如果发布的结果显示某个职业的患病率显著偏高，而张三的职业是公开信息，那么任何人都可以推断张三的风险偏高——这个推断在张三退出数据集之后依然成立，所以它不违反定义。差分隐私保护的是``你的参与没有让你更容易被推断''，不是``关于你的一切都推断不出来''。后者在数学上不可能：任何有用的统计发布都会改变对个体的推断，这正是它有用的原因。承诺后者的系统要么没在发布信息，要么没说实话。

\textbf{它也不管数据是怎么收集来的。}机制的随机性只覆盖从数据集到输出这一段。如果采集环节本身就有偏、或者原始数据被明文存放在某处，差分隐私一个字都没说。这和\hyperref[sec:13-Machine-Learning/01-Learning-Problems-and-Evaluation/04]{``数据划分、泄漏与分布偏移''}里那条判断同构：一个数学保证只覆盖它显式建模的那一段流程，流程之外的环节不会因为公式漂亮而自动安全。

到这里，隐私从一个含混的形容词变成了一个可以检验的不等式。剩下的问题是怎么造出满足它的机制。定义要求两个输出分布处处接近，而数据一变、真实答案就变——需要有东西把这个变化盖住。下一节说明该加多少噪声，以及决定这个量的既不是数据规模也不是 \(\varepsilon\) 本身，而是一个叫敏感度的量。
